\documentclass[11pt]{article}

\usepackage[T1]{fontenc}
\usepackage[utf8]{inputenc}
\usepackage{lmodern}
\usepackage[english]{babel}
\usepackage{amsmath,amssymb,amsthm}
\usepackage{microtype}
\usepackage{booktabs}
\usepackage{array}
\usepackage{geometry}
\usepackage{xcolor}
\usepackage{hyperref}

\geometry{left=2.8cm, right=2.8cm, top=2.8cm, bottom=2.8cm}
\setlength{\emergencystretch}{2em}

\definecolor{LINK}{rgb}{0.30,0.12,0.08}
\definecolor{URL}{rgb}{0.45,0.12,0.48}
\definecolor{CITE}{rgb}{0.03,0.28,0.08}
\hypersetup{
  colorlinks=true,
  linkcolor=LINK,
  urlcolor=URL,
  citecolor=CITE,
  pdfborder={0 0 0},
  linktocpage=true,
  pdftitle={Admissible qx+1 Sequences, Semiconvergents, and Rational Catalan Numbers},
  pdfauthor={Mike Winkler},
  pdfsubject={Admissible qx+1 sequences and extremal Beatty passage counts},
  pdfkeywords={qx+1 problem, Beatty sequence, semiconvergent, rational Catalan number, Dyck path, OEIS}
}

\newtheorem{theorem}{Theorem}
\newtheorem{proposition}[theorem]{Proposition}
\newtheorem{corollary}[theorem]{Corollary}
\newtheorem{lemma}[theorem]{Lemma}
\newtheorem{remark}[theorem]{Remark}
\newtheorem{definition}[theorem]{Definition}

\newcommand{\OEIS}[1]{\href{https://oeis.org/#1}{#1}}
\newcommand{\Aq}{\mathcal A}
\newcommand{\Lq}{\mathcal L}
\newcommand{\Uq}{\mathcal U}
\newcommand{\Cat}{\operatorname{Cat}}

\title{Admissible \(qx+1\) Sequences, Semiconvergents,\\
and Rational Catalan Numbers\\[5mm]}
\author{Mike Winkler\\[5mm]
Faculty of Mathematics, Ruhr University Bochum, Germany
\\mike.winkler@ruhr-uni-bochum.de}
\date{September 14, 2026\\[5mm]}

\begin{document}

\maketitle

\begin{abstract}
For an odd integer \(q\geq3\), let \(a_q(r)\) count finite words in
\(\{q/2,1/2\}\) containing exactly \(r\) factors \(q/2\), with every
proper prefix product greater than \(1\) and total product less than \(1\).
For \(q=3,5,7\), these are OEIS \OEIS{A100982}, \OEIS{A174795}, and
\OEIS{A174796}. Put
\[
 \alpha_q=\log_2q,\qquad m_r(q)=\lfloor r\alpha_q\rfloor.
\]
Using the classical cycle lemma in this \(q\)-specific setting, we recover
the two natural Beatty passage bounds
\[
 \frac1r\binom{m_r(q)-1}{r-1}
 \leq a_q(r)\leq
 \frac1r\binom{m_r(q)}{r-1},
\]
and prove their equality criteria directly. The lower and upper equality
orders are the denominators of the strict lower and upper one-sided best
approximations to \(\log_2q\); together they are the denominators of all
convergents and semiconvergents. Equivalently, they are the strict record
minima and maxima of the binary mantissas \(q^r/2^{m_r(q)}\). At every
nontrivial equality order, the admissible words are in explicit bijection
with rational Dyck paths, so that \(a_q(r)\) is a rational Catalan number.
We also prove monotonicity in \(q\), compare the three OEIS sequences, and
derive their growth constants and exact normalized oscillations.
\end{abstract}

\medskip
\noindent\textit{Keywords.}
\(qx+1\) problem, admissible sequence, Beatty sequence, continued fraction,
semiconvergent, rational Catalan number, Dyck path, OEIS.

\smallskip
\noindent\textit{2020 Mathematics Subject Classification.}
11B83, 11A55, 05A15, 11B37.

\section{Introduction}
\label{sec:introduction}

The integer sequence \OEIS{A100982} counts admissible multiplicative words
built from the two factors \(3/2\) and \(1/2\). It occurs naturally in the
coefficient dynamics of the \(3x+1\) map and was recorded in connection
with Wagon's treatment of the problem \cite{Wagon85,OEIS100982}. The
analogous \(5x+1\) and \(7x+1\) counts are \OEIS{A174795} and
\OEIS{A174796} \cite{OEIS174795,OEIS174796}. The OEIS entries record several recurrences, including formulas of
Zarubin; \OEIS{A100982} also contains formulas explicitly labelled
conjectural. Our purpose is not to replace those recurrences but to identify
the arithmetic structure of the orders at which the counts attain their
natural cycle bounds.

For odd \(q\geq3\), let \(a_q(r)\) denote the number of words containing
exactly \(r\) factors \(q/2\), such that the running product stays above
\(1\) until the final step and then crosses below \(1\). These words are
lower first-passage compositions relative to the Beatty boundary
\[
 m_i(q)=\lfloor i\log_2q\rfloor,
\]
which places \OEIS{A100982}, \OEIS{A174795}, and \OEIS{A174796} in a
single family.

General Beatty passage counts satisfy cycle bounds whose equality cases are
governed by one-sided Diophantine approximation; a broader marked-rotation
and lattice-factorization treatment is given in \cite{WinBeatty}. For the present paper, however, the required bounds and equality criteria
are reproved directly from the classical cycle lemma of Dvoretzky and
Motzkin \cite{DvoretzkyMotzkin47}; see also Dershowitz and Zaks
\cite{DershowitzZaks90}. This removes any logical dependence on the general
theory. When \(\alpha_q=\log_2q\), the
result has a particularly concrete form. If
\[
 M_q(r)=\frac{q^r}{2^{\lfloor r\log_2q\rfloor}},
 \qquad 1<M_q(r)<2,
\]
then the lower equality orders are exactly the strict record minima of
\(M_q(r)\), whereas the upper equality orders are exactly the strict record
maxima. Thus extremal values occur precisely when a power \(q^r\) gives a
new best multiplicative approach to a power of two from one of the two
sides.

Han\v{c}l and Turek proved that best one-sided approximations of the first
and second kinds coincide and described them by convergents and
semiconvergents \cite[Theorem~4.5]{HanclTurek}. Consequently the union of
the two equality sets is exactly the set of denominators of the convergents
and semiconvergents of \(\log_2q\). For \(q=3\), this is the set of
terms of OEIS \OEIS{A206788}; that entry lists the denominator \(1\)
twice, once for each side, while as an equality-order set it occurs only
once.

At equality orders there is also a direct enumerative interpretation. The
admissible words become rational Dyck paths in a coprime rectangle, and the
extremal values are rational Catalan numbers. For example,
\[
 A100982(7)=\Cat(7,4)=30,\qquad
 A174795(4)=\Cat(4,5)=14.
\]
For \OEIS{A174796}, the values \(1,2,7,30,143,728\) agree initially with
\OEIS{A006013}; Section~\ref{sec:catalan} explains exactly why this
coincidence persists through order \(6\) and then breaks.

For \(q=3\), a Beatty--Ferrers model with cut refinements, matroid structure,
and a power-of-two law was developed in \cite{WinFerrers}. The general
marked-rotation theory is the subject of \cite{WinBeatty}. The present paper
is intentionally narrower: its contribution is the \(qx+1\) sequence
synthesis, the self-contained extremal-order theorem, its semiconvergent and
power-of-two interpretations, the rational Dyck-path specialization at those
extremal orders, and the comparison across \(q\).

\section{Admissible words and Beatty passage counts}
\label{sec:words}

Fix an odd integer \(q\geq3\), and put
\[
 \alpha_q=\log_2q,\qquad
 m_r=m_r(q)=\lfloor r\alpha_q\rfloor.
\]
Since \(q\) is not a power of \(2\), \(\alpha_q\) is irrational.

\begin{definition}
A \emph{\(q\)-admissible word of order \(r\)} is a finite word
\[
 w=(x_1,\ldots,x_L),
 \qquad
 x_j\in\left\{\frac q2,\frac12\right\},
\]
containing exactly \(r\) occurrences of \(q/2\), such that
\[
 \prod_{j=1}^{\ell}x_j>1
 \quad(1\leq\ell<L),
 \qquad
 \prod_{j=1}^{L}x_j<1.
\]
Let \(\Aq_q(r)\) be the set of these words and set
\[
 a_q(r)=|\Aq_q(r)|.
\]
\end{definition}

All three OEIS sequences used here have offset \(1\), and our order \(r\)
uses that same index:
\[
 a_3(r)=A100982(r),\qquad
 a_5(r)=A174795(r),\qquad
 a_7(r)=A174796(r).
\]
The formula fields of the three OEIS entries write the prefix condition for
\(1<m<L\), omitting \(m=1\). Their example fields settle the intended
convention: in all three entries the unique word of order \(1\) begins
with \(q/2\), whereas the literal formula would also admit a word beginning
with \(1/2\). Thus Definition~1 agrees with the OEIS data and examples
\cite{OEIS100982,OEIS174795,OEIS174796}.

For later comparison across slopes, define for every real \(\alpha>1\)
\[
 c_1(\alpha)=1,
\]
and, for \(r\geq2\),
\begin{equation}
 c_r(\alpha)=
 \#\left\{
 1\leq s_1<\cdots<s_{r-1}:
 s_i\leq\lfloor i\alpha\rfloor\ (1\leq i<r)
 \right\}.
 \tag{2.1}
\end{equation}
This is the lower Beatty passage count in increasing-sequence form.

\begin{lemma}[Beatty passage representation]
\label{lem:beatty-passage}
Every word in \(\Aq_q(r)\) has length \(m_r+1\) and has a unique
representation
\[
 \frac q2\left(\frac12\right)^{b_1-1}
 \frac q2\left(\frac12\right)^{b_2-1}
 \cdots
 \frac q2\left(\frac12\right)^{b_r-1},
\]
where \(b_i\geq1\). If \(B_i=b_1+\cdots+b_i\), then the word is
admissible if and only if
\[
 B_r=m_r+1,
 \qquad
 B_i\leq m_i
 \quad(1\leq i<r).
 \tag{2.2}
\]
Consequently
\[
 a_q(r)=c_r(\alpha_q).
 \tag{2.3}
\]
\end{lemma}

\begin{proof}
The last letter must be \(1/2\), since multiplication by \(q/2>1\) cannot
produce the first passage below \(1\). Hence
\[
 \frac{q^r}{2^{L-1}}>1,
 \qquad
 \frac{q^r}{2^L}<1,
\]
so \(2^{L-1}<q^r<2^L\). Because \(r\alpha_q\notin\mathbb Z\), this gives
\(L=m_r+1\).

The first letter must be \(q/2\). Grouping each occurrence of \(q/2\)
with the following string of factors \(1/2\) gives the stated block
decomposition and \(B_r=L\). At the end of block \(i\) the running product
is \(q^i/2^{B_i}\). After the initial factor of a block the product decreases,
so the block endpoint is its smallest value. In the last block every proper
prefix is still above \(1\), because the product immediately before the
last letter is above \(1\). Hence only the first \(r-1\) block endpoints
need be tested. Since \(i\alpha_q\) is nonintegral,
\[
 \frac{q^i}{2^{B_i}}>1
 \iff B_i<i\alpha_q
 \iff B_i\leq m_i.
\]
Thus the admissible words are in bijection with the increasing partial sums
counted in (2.1). The final part is automatically positive because
\(B_{r-1}\leq m_{r-1}<m_r+1\).
\end{proof}

\section{Cycle bounds and extremal orders}
\label{sec:extremal}

We first record the cyclic-rotation fact used in both bounds.  It is the
standard partial-sum form of the classical cycle lemma of Dvoretzky and
Motzkin \cite{DvoretzkyMotzkin47}; a modern account and applications are
given by Dershowitz and Zaks \cite{DershowitzZaks90}.

\begin{lemma}[Cyclic rotation]
\label{lem:rotation}
Let \(x=(x_1,\ldots,x_r)\) be a positive composition of an integer \(T\),
and let \(X_i=x_1+\cdots+x_i\). Every cyclic orbit contains a rotation
satisfying
\[
 X_i\leq\frac{iT}{r}
 \qquad(1\leq i<r).
 \tag{3.1}
\]
If \(\gcd(r,T)=1\), every orbit has length \(r\) and contains exactly one
rotation satisfying the strict inequalities \(X_i<iT/r\) for all
\(i<r\). For arbitrary \(r,T\), an orbit containing such a strict rotation
has full length and contains at most one of them.
\end{lemma}

\begin{proof}
Put \(S_0=0\) and
\[
 S_i=rX_i-iT\qquad(1\leq i\leq r).
\]
Starting immediately after a maximum of
\(S_0,\ldots,S_{r-1}\) gives a rotation for which all proper transformed
partial sums are nonpositive, proving (3.1). If \(\gcd(r,T)=1\), the values
\(S_0,\ldots,S_{r-1}\) are pairwise distinct modulo \(r\), hence the
maximum is unique and no proper transformed sum vanishes. This gives the
unique strict rotation. Conversely a shorter period forces equality with
the rational line after one period, so a strict rotation can occur only in
a full orbit. Extending \(S_i\) periodically by \(S_{i+r}=S_i\), a strict
rotation beginning after index \(k\) means
\[
 S_i<S_k\qquad(k<i<k+r).
\]
Thus \(S_k\) is the strict maximum on one complete residue system modulo
\(r\), and such a starting point is unique.
\end{proof}

For an irrational \(\alpha>1\), write
\[
 m_i=\lfloor i\alpha\rfloor,
 \qquad
 \delta_i=\{i\alpha\}.
\]
The next elementary lemma converts the record conditions into rational
floor conditions.

\begin{lemma}[Record--floor equivalence]
\label{lem:record-floor}
Let \(r\geq2\), \(m=m_r\), and \(t=m+1\). Then
\[
 \delta_r<\delta_j\quad(1\leq j<r)
 \tag{3.2}
\]
if and only if
\[
 \gcd(r,m)=1,
 \qquad
 m_j=\left\lfloor\frac{mj}{r}\right\rfloor
 \quad(1\leq j<r).
 \tag{3.3}
\]
Similarly,
\[
 \delta_r>\delta_j\quad(1\leq j<r)
 \tag{3.4}
\]
if and only if
\[
 \gcd(r,t)=1,
 \qquad
 m_j=\left\lfloor\frac{tj}{r}\right\rfloor
 \quad(1\leq j<r).
 \tag{3.5}
\]
\end{lemma}

\begin{proof}
Assume first (3.2). If \(d>1\) divides both \(r\) and \(m\), then with
\(j=r/d\) one has \(\delta_j=\delta_r/d<\delta_r\), a contradiction.
Thus \(\gcd(r,m)=1\). Since \(m/r<\alpha\),
\(\lfloor mj/r\rfloor\leq m_j\). If the inequality were strict for some
\(j<r\), then
\[
 \delta_j=\alpha j-m_j
 <\alpha j-\frac{mj}{r}
 =\frac{j}{r}\delta_r<\delta_r,
\]
again a contradiction. Hence (3.3) holds.

Conversely assume (3.3), and suppose \(\delta_j<\delta_r\) for some
\(j<r\). Since
\(\delta_j+\delta_{r-j}\in\{\delta_r,\delta_r+1\}\), the second
alternative is impossible, so
\(\delta_{r-j}=\delta_r-\delta_j\). The two floor equalities in (3.3)
imply
\[
 \delta_j\geq\frac{j}{r}\delta_r,
 \qquad
 \delta_{r-j}\geq\frac{r-j}{r}\delta_r.
\]
Substitution gives the reverse inequality for \(\delta_j\), hence equality.
Then \(m_j=mj/r\), so \(r\mid mj\). Since \(\gcd(r,m)=1\), this forces
\(r\mid j\), impossible for \(1\leq j<r\). Thus (3.2) holds.

For the upper statement, assume (3.4). If \(d>1\) divides \(r\) and
\(t\), then with \(j=r/d\), the integer \(t/d\) is the least integer
above \(\alpha j\), and
\[
 1-\delta_j=\frac{1-\delta_r}{d}<1-\delta_r,
\]
contradicting (3.4). Hence \(\gcd(r,t)=1\). Since \(t/r>\alpha\), if
\(\lfloor tj/r\rfloor>m_j\), then
\[
 1-\delta_j
 \leq\frac{tj}{r}-\alpha j
 =\frac{j}{r}(1-\delta_r)<1-\delta_r,
\]
again a contradiction. This proves (3.5).

Conversely assume (3.5) and suppose \(\delta_j>\delta_r\). Then
\(\delta_j+\delta_{r-j}=\delta_r+1\), so
\[
 (1-\delta_j)+(1-\delta_{r-j})=1-\delta_r.
\]
The two floor equalities in (3.5) imply
\[
 1-\delta_j>\frac{j}{r}(1-\delta_r),
 \qquad
 1-\delta_{r-j}>\frac{r-j}{r}(1-\delta_r),
\]
whose sum is impossible. Thus (3.4) follows.
\end{proof}

\begin{remark}[Mechanical-word interpretation]
The floor conditions in Lemma~\ref{lem:record-floor} are classical in
the theory of mechanical, Sturmian, and Christoffel words: the irrational
mechanical prefix agrees with the corresponding rational prefix of slope
\(m_r/r\) or \((m_r+1)/r\). Its relation with convergents and
semiconvergents is standard; see Berstel and S\'e\'ebold
\cite{BerstelSeebold02}. The lemma is included only to keep the
extremal-count argument self-contained.
\end{remark}

Return now to \(\alpha_q=\log_2q\), and abbreviate \(m_r=m_r(q)\). Put
\[
 C_q(r)=\binom{m_r-1}{r-1},
 \qquad
 D_q(r)=\binom{m_r}{r-1},
 \qquad
 M_q(r)=\frac{q^r}{2^{m_r}}=2^{\delta_r}.
 \tag{3.6}
\]
Since \(\alpha_q\geq\log_2 3>3/2\), one has
\[
 m_r\geq r\quad(r\geq1),
 \qquad
 m_r\geq r+1\quad(r\geq2).
 \tag{3.7}
\]

\begin{theorem}[Extremal orders]
\label{thm:extremal-orders}
For every odd \(q\geq3\) and every \(r\geq1\),
\[
 \frac{C_q(r)}{r}
 \leq a_q(r)\leq
 \frac{D_q(r)}{r}.
 \tag{3.8}
\]
For \(r\geq2\), lower equality holds if and only if \(\delta_r\) is a
strict record minimum,
\[
 a_q(r)=\frac{C_q(r)}{r}
 \iff
 \delta_r<\delta_j\quad(1\leq j<r),
 \tag{3.9}
\]
and upper equality holds if and only if \(\delta_r\) is a strict record
maximum,
\[
 a_q(r)=\frac{D_q(r)}{r}
 \iff
 \delta_r>\delta_j\quad(1\leq j<r).
 \tag{3.10}
\]
Equivalently, the lower and upper equality orders are respectively the
strict record minima and maxima of \(M_q(r)\). The index \(r=1\) is the
unique common equality order.
\end{theorem}

\begin{proof}
Let \(m=m_r\) and \(t=m+1\).

For the lower bound, consider all positive compositions \(h\) of \(m\)
into \(r\) parts. From each cyclic orbit choose one rotation supplied by
Lemma~\ref{lem:rotation}, with partial sums
\(H_i\leq im/r<\alpha_q i\), hence \(H_i\leq m_i\). Adding \(1\) to
its last part produces an admissible composition of total \(t\). This
choice defines an injection from cyclic orbits to admissible words: if two
chosen images were equal, subtracting \(1\) from the last part would give
the same rotated composition and hence the same orbit. Since there are
\(C_q(r)\) positive compositions of \(m\) and every orbit has at most
\(r\) elements,
\[
 a_q(r)\geq C_q(r)/r.
\]

If \(d=\gcd(r,m)>1\), a positive composition of \(m/d\) into \(r/d\)
parts repeated \(d\) times gives a short orbit. Hence there are strictly
more than \(C_q(r)/r\) orbits, so the lower inequality is strict. Assume
therefore \(\gcd(r,m)=1\). Every orbit then has length \(r\) and exactly
one rotation strictly below the rational line \(im/r\). If
\[
 m_i=\left\lfloor\frac{mi}{r}\right\rfloor
 \quad(1\leq i<r),
 \tag{3.11}
\]
then (3.11) at \(i=r-1\) gives \(m_{r-1}\leq m-2\), and every
admissible composition satisfies
\[
 b_r=t-B_{r-1}\geq t-m_{r-1}\geq3.
\]
Subtracting \(1\) from its last part gives a composition \(h\) of \(m\)
with
\[
 H_i=B_i\leq m_i=
 \left\lfloor\frac{mi}{r}\right\rfloor<\frac{mi}{r}
 \qquad(1\leq i<r),
\]
where the last inequality uses \(\gcd(r,m)=1\). Thus every admissible
word maps to the unique strict rational rotation of its orbit, so there is
at most one admissible word per orbit. Together with the lower bound this
gives \(a_q(r)=C_q(r)/r\).

If (3.11) fails, the composition of \(m\) with partial sums \(H_i=m_i\)
for \(i<r\) and \(H_r=m\) has positive parts because
\(m_i-m_{i-1}\geq1\). Adding \(1\) to its last part gives one
Beatty-admissible word, while the unique strict rational rotation in the
same orbit gives another after the same addition. They are distinct
because the boundary composition lies above the rational line at a
discrepancy index. Hence one orbit contributes at least two words, whereas
every orbit contributes at least one, so the lower inequality is strict.
By Lemma~\ref{lem:record-floor}, this proves (3.9).

For the upper bound, every admissible composition \(b\) of \(t\) satisfies
\[
 B_i\leq m_i<\alpha_q i<\frac{ti}{r}
 \qquad(1\leq i<r).
\]
Thus it is a strict rationally-below rotation. By Lemma~\ref{lem:rotation},
a contributing orbit has full length and contributes at most one word.
There are \(D_q(r)=\binom{t-1}{r-1}\) positive compositions of \(t\), so
\(a_q(r)\leq D_q(r)/r\).

If \(\gcd(r,t)>1\), a short orbit exists; by the final assertion of
Lemma~\ref{lem:rotation}, no member of such an orbit can be a strict
rationally-below rotation. Hence that orbit contributes nothing and the
upper inequality is strict. If \(\gcd(r,t)=1\), all orbits have length
\(r\) and exactly one strict rationally-below rotation. Equality holds
precisely when the Beatty and rational floor boundaries agree,
\[
 m_i=\left\lfloor\frac{ti}{r}\right\rfloor
 \quad(1\leq i<r).
\]
Indeed, if a discrepancy occurs, the rational boundary composition with
partial sums \(\lfloor ti/r\rfloor\) has positive parts because \(t/r>1\).
It is the unique strict rational rotation
of its orbit but violates the Beatty condition at that index, so the upper
inequality is strict. Lemma~\ref{lem:record-floor} now gives (3.10).

Finally \(M_q(r)=2^{\delta_r}\), so its record minima and maxima are exactly
those of \(\delta_r\). For \(r\geq2\), the two equality values differ
because
\[
 \frac{D_q(r)}{C_q(r)}=
 \frac{m_r}{m_r-r+1}>1.
\]
\end{proof}

The theorem has a direct multiplicative interpretation. At a lower equality
order, \(q^r\) gives a new smallest ratio to its lower bounding power
\(2^{m_r}\). At an upper equality order,
\[
 \frac{2^{m_r+1}}{q^r}=\frac{2}{M_q(r)}
\]
gives a new smallest ratio to the next power of two.

\section{Continued fractions and semiconvergents}
\label{sec:cf}

Write
\[
 \alpha_q=[a_0;a_1,a_2,\ldots],
\]
and let \(P_k/Q_k\) be the \(k\)-th convergent, with
\[
 P_{-1}=1,\quad P_0=a_0,
 \qquad
 Q_{-1}=0,\quad Q_0=1.
\]
For an irrational \(\alpha>1\), the rational number
\(\lfloor r\alpha\rfloor/r<\alpha\) is called a \emph{strict best
lower one-sided approximation of the second kind} if
\[
 \{r\alpha\}<\{j\alpha\}
 \qquad(1\leq j<r),
\]
and \(\lceil r\alpha\rceil/r>\alpha\) is a \emph{strict best upper
one-sided approximation of the second kind} if
\[
 1-\{r\alpha\}<1-\{j\alpha\}
 \qquad(1\leq j<r).
\]
Thus the record conditions in Theorem~\ref{thm:extremal-orders} are
precisely the lower and upper one-sided best-approximation conditions of
the second kind. Earlier descriptions of best lower and upper approximates
in terms of principal and intermediate convergents are due to Kimberling
\cite{Kimberling97}. Han\v{c}l and Turek prove that the first- and
second-kind one-sided notions coincide and give the parametrization used
below \cite[Theorem~4.5]{HanclTurek}.

\begin{theorem}[Continued-fraction parametrization]
\label{thm:cf}
The lower equality orders are
\[
 \Lq_q=
 \left\{
 sQ_k+Q_{k-1}:
 k\geq1\ \text{odd},\
 0\leq s<a_{k+1}
 \right\},
 \tag{4.1}
\]
and the upper equality orders are
\[
 \Uq_q=
 \left\{
 sQ_k+Q_{k-1}:
 k\geq0\ \text{even},\
 0\leq s<a_{k+1},\
 (k,s)\neq(0,0)
 \right\}.
 \tag{4.2}
\]
Their union is exactly the set of denominators of all convergents and
semiconvergents of \(\log_2q\). The two sets meet only at \(r=1\).
\end{theorem}

\begin{proof}[Derivation from Han\v{c}l--Turek]
Writing their convergent and semiconvergent parametrization in our indexing
\(P_{-1}=1,P_0=a_0,Q_{-1}=0,Q_0=1\), the one-sided best approximations are
\[
 \frac{sP_k+P_{k-1}}{sQ_k+Q_{k-1}},
 \qquad 0\leq s<a_{k+1},
\]
with odd \(k\) on the lower side and even \(k\) on the upper side. The
pair \((k,s)=(0,0)\) is excluded because its denominator is
\(Q_{-1}=0\). Taking denominators gives (4.1) and (4.2). The union statement
is the same theorem. Theorem~\ref{thm:extremal-orders} shows that a positive
denominator can belong to both equality sets only for \(r=1\). The lower
set always contains \(1\) through \((k,s)=(1,0)\). On the upper side, if
\(a_1\geq2\), the index \(1\) comes from \((k,s)=(0,1)\); if
\(a_1=1\), then \(Q_1=Q_0=1\) and it comes from \((k,s)=(2,0)\).
Thus \(1\) occurs once on each side.
\end{proof}

In particular, the principal convergents already show that both equality
sets are infinite: with the above indexing, the even convergents
\(P_{2n}/Q_{2n}\) lie below \(\alpha_q\) and give lower record
denominators, whereas the odd convergents \(P_{2n+1}/Q_{2n+1}\) lie
above \(\alpha_q\) and give upper record denominators.

For \(q=3\),
\[
 \log_2 3=[1;1,1,2,2,3,1,5,2,23,\ldots].
\]
The two equality sets begin
\[
 \Lq_3:1,2,7,12,53,359,665,\ldots,
 \qquad
 \Uq_3:1,3,5,17,29,41,94,147,200,253,306,971,\ldots.
\]
The union \(\Lq_3\cup\Uq_3\) is exactly the set of terms of OEIS
\OEIS{A206788}, the denominators of semiconvergents to \(\log_2 3\)
\cite{OEIS206788}. The OEIS entry itself lists the denominator \(1\)
twice, once for each side; as an equality-order set it occurs only once.

For comparison, Table~\ref{tab:records} lists the equality orders up to
\(20\) for the three OEIS cases.

\begin{table}[ht]
\centering
\small
\begin{tabular}{c p{5.3cm} p{5.3cm}}
\toprule
\(q\) & lower equality orders \(r\leq20\) & upper equality orders \(r\leq20\)\\
\midrule
3 & \(1,2,7,12\) & \(1,3,5,17\)\\
5 & \(1,4,7,10,13,16,19\) & \(1,2,3\)\\
7 & \(1,2,3,4,5\) & \(1,6,11,16\)\\
\bottomrule
\end{tabular}
\caption{Initial lower and upper equality orders.}
\label{tab:records}
\end{table}

For \(q=3\), several auxiliary sequences are also catalogued:
\(m_r(3)\) is \OEIS{A056576}, the word length \(m_r(3)+1\) is
\OEIS{A020914}, the continued fraction of \(\log_2 3\) is
\OEIS{A028507}, and the principal-convergent denominators are
\OEIS{A005664}.

The full merged denominator sequences for \(q=5\) and \(q=7\) are
recorded in Section~\ref{sec:conclusion}.

\section{Rational Catalan values}
\label{sec:catalan}

For coprime positive integers \(a,b\), write
\[
 \Cat(a,b)=\frac1{a+b}\binom{a+b}{a}.
 \tag{5.1}
\]
This is the rational Catalan number counting lattice paths from
\((0,0)\) to \((a,b)\) that stay weakly below the diagonal; see Bizley
\cite{Bizley54} and Armstrong, Rhoades, and Williams \cite{ArmstrongRW13}.
The OEIS entry for \OEIS{A100982} already records lattice-walk material,
including a link by van Tol and cross-references to \OEIS{A060941} and
\OEIS{A293946} \cite{OEIS100982}. The point here is different: the
extremal-order theorem identifies exactly when the irrational Beatty
boundary collapses to a coprime rational boundary, and does so uniformly
for the \(q\)-family.

\begin{theorem}[Dyck paths at extremal orders]
\label{thm:dyck}
Let \(r\geq2\). By (3.7), all path heights below are positive, so the
rectangles are nondegenerate. If \(r\in\Lq_q\), then
\[
 a_q(r)=\Cat(r,m_r-r)
 =\frac1{m_r}\binom{m_r}{r}.
 \tag{5.2}
\]
More precisely, \(\Aq_q(r)\) is in bijection with the rational Dyck paths
from \((0,0)\) to \((r,m_r-r)\).

If \(r\in\Uq_q\), then
\[
 a_q(r)=\Cat(r,m_r+1-r)
 =\frac1{m_r+1}\binom{m_r+1}{r}.
 \tag{5.3}
\]
In this case \(\Aq_q(r)\) is in bijection with the rational Dyck paths
from \((0,0)\) to \((r,m_r+1-r)\).
\end{theorem}

\begin{proof}
Suppose first that \(r\in\Lq_q\). By Lemma~\ref{lem:record-floor},
\[
 m_i=\left\lfloor\frac{m_ri}{r}\right\rfloor
 \quad(1\leq i<r),
 \qquad
 \gcd(r,m_r)=1.
 \tag{5.4}
\]
Let \((b_1,\ldots,b_r)\) be the composition associated with an admissible
word. Since \(\alpha_q>1\), one has \(m_r-m_{r-1}\geq1\), and hence
\[
 b_r=m_r+1-B_{r-1}
 \geq m_r+1-m_{r-1}\geq2.
\]
Subtracting \(1\) from the last part therefore gives, reversibly, a positive
composition \(h\) of \(m_r\). This subtraction is needed only on the lower
side because the admissible word has total \(m_r+1\), whereas the lower
rational boundary in (5.4) has total \(m_r\).

Set \(y_i=h_i-1\) and \(s=m_r-r\). Then \(y_i\geq0\),
\(\sum_i y_i=s\), and
\[
 Y_i:=y_1+\cdots+y_i
 \leq
 \left\lfloor\frac{m_ri}{r}\right\rfloor-i
 =\left\lfloor\frac{si}{r}\right\rfloor.
\]
Thus
\[
 (y_1,\ldots,y_r)
 \longmapsto
 E N^{y_1}E N^{y_2}\cdots E N^{y_r}
\]
is a bijection onto the paths from \((0,0)\) to \((r,s)\) lying weakly
below the diagonal. Since \(s=m_r-r\),
\[
 \gcd(r,s)=\gcd(r,m_r)=1,
\]
and their number is (5.2).

If \(r\in\Uq_q\), put \(t=m_r+1\). Lemma~\ref{lem:record-floor} gives
\[
 m_i=\left\lfloor\frac{ti}{r}\right\rfloor
 \quad(1\leq i<r),
 \qquad
 \gcd(r,t)=1.
 \tag{5.5}
\]
Now no subtraction is needed because the admissible composition already
has total \(t\). Set \(y_i=b_i-1\) and \(s=t-r\). Then
\[
 Y_i=B_i-i
 \leq
 \left\lfloor\frac{ti}{r}\right\rfloor-i
 =\left\lfloor\frac{si}{r}\right\rfloor,
\]
and the same path construction gives (5.3), because
\[
 \gcd(r,s)=\gcd(r,t-r)=\gcd(r,t)=1.
\]
\end{proof}

The theorem explains several small values without recursion. For
\(\OEIS{A100982}\),
\[
 A100982(7)=\Cat(7,4)=30,\qquad
 A100982(12)=\Cat(12,7)=2652,
\]
and the upper extremal order \(17\) gives
\(A100982(17)=\Cat(17,10)=312455\).

For \(q=5\), the first three nontrivial values are ordinary Catalan numbers:
\[
 A174795(2)=2,\qquad A174795(3)=5,\qquad A174795(4)=14,
\]
because in these cases Theorem~\ref{thm:dyck} gives \(\Cat(r,r+1)\).
This coincidence breaks at \(r=5\): \(A174795(5)=56\), whereas
\(\OEIS{A000108}(5)=42\).

For \(q=7\),
\[
 A174796(1),\ldots,A174796(6)=1,2,7,30,143,728,
\]
which is exactly \(\OEIS{A006013}(0),\ldots,\OEIS{A006013}(5)\). For \(2\leq r\leq5\)
one has \(m_r=3r-1\) and a lower equality order, while at \(r=6\) one has
\(m_6+1=17=3\cdot6-1\) and an upper equality order. Hence all five nontrivial
terms equal
\[
 \Cat(r,2r-1)
 =\frac1r\binom{3r-2}{r-1}
 =\OEIS{A006013}(r-1).
\]
The coincidence breaks at \(r=7\), where \(A174796(7)=3148\) but
\(\OEIS{A006013}(6)=3876\).

\section{Comparison across \texorpdfstring{\(q\)}{q}}
\label{sec:comparison}

The Beatty representation also gives a sharp monotonicity statement.

\begin{proposition}[Monotonicity and rigidity]
\label{prop:monotonicity}
Let \(1<\alpha<\beta\) and \(r\geq2\).  Then
\[
 c_r(\alpha)\leq c_r(\beta).
 \tag{6.1}
\]
Moreover,
\[
 c_r(\alpha)=c_r(\beta)
 \iff
 \lfloor i\alpha\rfloor=\lfloor i\beta\rfloor
 \quad(1\leq i<r).
 \tag{6.2}
\]
Hence, for fixed \(r\), the function
\(\alpha\mapsto c_r(\alpha)\) is a nondecreasing step function whose
constant cells are exactly the cells of the finite Beatty prefix
\[
 \bigl(\lfloor\alpha\rfloor,\ldots,
       \lfloor(r-1)\alpha\rfloor\bigr).
\]
\end{proposition}

\begin{proof}
Write
\[
 M_i=\lfloor i\alpha\rfloor,\qquad
 N_i=\lfloor i\beta\rfloor.
\]
By definition (2.1), \(c_r(\alpha)\) counts
\[
 1\leq s_1<\cdots<s_{r-1},\qquad s_i\leq M_i,
\]
and the same description with \(N_i\) counts \(c_r(\beta)\).
Since \(M_i\leq N_i\), (6.1) follows.

If all \(M_i=N_i\), equality is immediate.  Conversely, assume
\(M_j<N_j\) for some \(j<r\).  Since \(\alpha>1\), \(M_j\geq j\).
Set
\[
 s_i=i\quad(i<j),\qquad
 s_j=M_j+1,
\]
and
\[
 s_i=M_j+1+i-j\quad(j<i<r).
\]
This sequence is strictly increasing and violates the \(\alpha\)-bound at
\(j\).  On the other hand \(s_j\leq N_j\), and for \(i>j\),
\[
 N_i-N_j\geq i-j
\]
because \(\beta>1\).  Hence \(s_i\leq N_i\), so the sequence is counted
for \(\beta\) but not for \(\alpha\).  Thus the inequality is strict.
\end{proof}

\begin{corollary}
\label{cor:three-sequences}
For every \(r\geq2\),
\[
 A100982(r)<A174795(r).
 \tag{6.3}
\]
Furthermore,
\[
 A174795(2)=A174796(2)=2,
\]
and
\[
 A174795(r)<A174796(r)
 \qquad(r\geq3).
 \tag{6.4}
\]
\end{corollary}

\begin{proof}
We have
\[
 \log_2 3<\log_2 5<\log_2 7.
\]
For the first pair the Beatty prefixes already differ at \(i=1\), since
\[
 \lfloor\log_2 3\rfloor=1,\qquad
 \lfloor\log_2 5\rfloor=2.
\]
For the second pair the \(i=1\) values are both \(2\), but
\[
 \lfloor2\log_2 5\rfloor=4,\qquad
 \lfloor2\log_2 7\rfloor=5.
\]
Now apply Proposition~\ref{prop:monotonicity}.
\end{proof}

\section{Growth and exact normalized oscillation}
\label{sec:growth}

Put
\[
 B(\alpha)=
 \frac{\alpha^\alpha}{(\alpha-1)^{\alpha-1}},
 \qquad
 B_q=B(\alpha_q).
 \tag{7.1}
\]

\begin{corollary}[Growth constants]
\label{cor:growth}
For every odd \(q\geq3\),
\[
 \lim_{r\to\infty}a_q(r)^{1/r}=B_q.
 \tag{7.2}
\]
The function \(q\mapsto B_q\) is strictly increasing. In particular,
\[
 \lim_{r\to\infty}
 \left(\frac{A174795(r)}{A100982(r)}\right)^{1/r}
 =\frac{B_5}{B_3}
 \approx1.7218785536,
 \tag{7.3}
\]
and
\[
 \lim_{r\to\infty}
 \left(\frac{A174796(r)}{A174795(r)}\right)^{1/r}
 =\frac{B_7}{B_5}
 \approx1.2726660717.
 \tag{7.4}
\]
\end{corollary}

\begin{proof}
By Theorem~\ref{thm:extremal-orders}, \(a_q(r)\) is squeezed between
\(C_q(r)/r\) and \(D_q(r)/r\). Since
\[
 m_r=\alpha_q r+O(1),
 \qquad
 m_r-r=(\alpha_q-1)r+O(1)\longrightarrow\infty,
\]
Stirling's formula gives
\[
 \frac1r\log\binom{m_r-1}{r-1}
 =
 \alpha_q\log\alpha_q
 -(\alpha_q-1)\log(\alpha_q-1)+o(1),
\]
and the same limit with \(m_r\) in place of \(m_r-1\).  Hence both
bounds have \(r\)-th-root limit
\[
 \frac{\alpha_q^{\alpha_q}}
 {(\alpha_q-1)^{\alpha_q-1}}=B_q.
\]
Moreover,
\[
 \frac{d}{d\alpha}\log B(\alpha)
 =\log\frac{\alpha}{\alpha-1}>0,
\]
so \(q\mapsto B_q\) is strictly increasing. Taking quotients of the root
limits gives (7.3) and (7.4).
\end{proof}

Define
\[
 R_q(r)=
 \frac{r\,a_q(r)}
 {\binom{m_r(q)-1}{r-1}}.
 \tag{7.5}
\]

\begin{corollary}[Exact normalized oscillation]
\label{cor:oscillation}
For every odd \(q\geq3\),
\[
 \liminf_{r\to\infty}R_q(r)=1,
 \qquad
 \limsup_{r\to\infty}R_q(r)
 =\frac{\alpha_q}{\alpha_q-1}.
 \tag{7.6}
\]
The value \(R_q(r)=1\) occurs exactly at \(r\in\Lq_q\). At
\(r\in\Uq_q\),
\[
 R_q(r)=\frac{m_r(q)}{m_r(q)-r+1},
 \tag{7.7}
\]
and these values tend to the limsup. In fact
\[
 R_q(r)<\frac{\alpha_q}{\alpha_q-1}
 \qquad(r\geq1).
 \tag{7.8}
\]
\end{corollary}

\begin{proof}
Theorem~\ref{thm:extremal-orders} gives
\[
 1\leq R_q(r)\leq\frac{m_r}{m_r-r+1}.
\]
As noted after Theorem~\ref{thm:cf}, the alternating principal
convergents already supply infinitely many lower and upper equality
orders. Along \(\Lq_q\) the left endpoint is attained, and along
\(\Uq_q\) the right endpoint is attained. Since \(m_r/r\to\alpha_q\),
the latter values tend to \(\alpha_q/(\alpha_q-1)\), proving (7.6).
Finally
\[
 \frac{m_r}{m_r-r+1}<\frac{\alpha_q}{\alpha_q-1}
\]
is equivalent to \(\{r\alpha_q\}<\alpha_q\), which is automatic because
\(0<\{r\alpha_q\}<1<\alpha_q\).
\end{proof}

For reference, Table~\ref{tab:threeq} gives the basic constants to ten
decimal places.

\begin{table}[ht]
\centering
\small
\renewcommand{\arraystretch}{1.15}
\begin{tabular}{c c c c}
\toprule
\(q\) & \(\alpha_q=\log_2q\) & \(B_q\) &
\(\alpha_q/(\alpha_q-1)\)\\
\midrule
3 & 1.5849625007 & 2.8395137305 & 2.7095112914\\
5 & 2.3219280949 & 4.8892977952 & 1.7564707974\\
7 & 2.8073549221 & 6.2224434185 & 1.5532947557\\
\bottomrule
\end{tabular}
\caption{Exponential growth constants and normalized upper envelopes.}
\label{tab:threeq}
\end{table}

\section{OEIS consequences and concluding remarks}
\label{sec:conclusion}

For \OEIS{A100982}, Theorem~\ref{thm:cf} identifies the equality-order
set with the set of terms of \OEIS{A206788}; that OEIS entry records the
initial denominator \(1\) twice, once for each side. For \(q=5\), the
side-labelled merged denominator
sequence begins
\[
 1,1,2,3,4,7,10,13,16,19,22,25,28,31,59,87,146,205,351,497,643,\ldots,
\]
and for \(q=7\) it begins
\[
 1,1,2,3,4,5,6,11,16,21,26,31,57,83,109,135,244,353,462,571,\ldots.
\]
In these side-labelled lists the repeated initial \(1\) records the two
sides; as an equality-order set it occurs only once. These two sequences
are natural companions to
\OEIS{A174795} and \OEIS{A174796}.

The resulting picture is compact. Extremal admissible-word counts are
simultaneously cycle-bound equalities, one-sided best approximations,
semiconvergent denominators, and record approaches of \(q^r\) to powers of
\(2\). At those extremal orders the words reduce canonically to rational
Dyck paths and the counts to rational Catalan numbers. Proposition~\ref{prop:monotonicity}
then supplies a complementary rigidity principle across the parameter
\(q\). Together these statements give a common arithmetic and combinatorial
framework for \OEIS{A100982}, \OEIS{A174795}, and \OEIS{A174796} without
requiring the broader marked-rotation machinery of \cite{WinBeatty}.

\begin{thebibliography}{99}

\bibitem{ArmstrongRW13}
D.~Armstrong, B.~Rhoades, and N.~Williams,
\newblock Rational associahedra and noncrossing partitions,
\newblock \emph{Electron. J. Combin.} \textbf{20} (2013), no.~3, Paper P54.
\newblock \href{https://doi.org/10.37236/3432}{doi:10.37236/3432}.

\bibitem{BerstelSeebold02}
J.~Berstel and P.~S\'e\'ebold,
\newblock Sturmian words,
\newblock in M.~Lothaire, \emph{Algebraic Combinatorics on Words},
Encyclopedia of Mathematics and its Applications, vol.~90,
Cambridge University Press, Cambridge, 2002, pp.~45--110.
\newblock \href{https://doi.org/10.1017/CBO9781107326019.003}{doi:10.1017/CBO9781107326019.003}.

\bibitem{Bizley54}
M.~T.~L. Bizley,
\newblock Derivation of a new formula for the number of minimal lattice
paths from \((0,0)\) to \((km,kn)\) having just \(t\) contacts with the
line \(my=nx\) and having no points above this line; and a proof of
Grossman's formula for the number of paths which may touch but do not rise
above this line,
\newblock \emph{J. Inst. Actuaries} \textbf{80} (1954), 55--62.
\newblock \href{https://doi.org/10.1017/S002026810005424X}{doi:10.1017/S002026810005424X}.

\bibitem{DershowitzZaks90}
N.~Dershowitz and S.~Zaks,
\newblock The cycle lemma and some applications,
\newblock \emph{European J. Combin.} \textbf{11} (1990), 35--40.
\newblock \href{https://doi.org/10.1016/S0195-6698(13)80053-4}{doi:10.1016/S0195-6698(13)80053-4}.

\bibitem{DvoretzkyMotzkin47}
A.~Dvoretzky and Th.~Motzkin,
\newblock A problem of arrangements,
\newblock \emph{Duke Math. J.} \textbf{14} (1947), 305--313.
\newblock \href{https://doi.org/10.1215/S0012-7094-47-01423-3}{doi:10.1215/S0012-7094-47-01423-3}.

\bibitem{HanclTurek}
J.~Han\v{c}l and O.~Turek,
\newblock One-sided Diophantine approximations,
\newblock \emph{J. Phys. A: Math. Theor.} \textbf{52} (2019), 045205.
\newblock \href{https://doi.org/10.1088/1751-8121/aaf5d3}{doi:10.1088/1751-8121/aaf5d3}.


\bibitem{Kimberling97}
C.~Kimberling,
\newblock Best lower and upper approximates to irrational numbers,
\newblock \emph{Elem. Math.} \textbf{52} (1997), 122--126.
\newblock \href{https://doi.org/10.1007/S000170050020}{doi:10.1007/S000170050020}.

\bibitem{OEIS100982}
OEIS Foundation Inc.,
\newblock \emph{The On-Line Encyclopedia of Integer Sequences},
\newblock \href{https://oeis.org/A100982}{A100982},
Number of admissible sequences of order \(n\); related to \(3x+1\)
problem and Wagon's constant.

\bibitem{OEIS174795}
OEIS Foundation Inc.,
\newblock \emph{The On-Line Encyclopedia of Integer Sequences},
\newblock \href{https://oeis.org/A174795}{A174795},
Number of admissible sequences of order \(j\); related to \(5x+1\)
problem.

\bibitem{OEIS174796}
OEIS Foundation Inc.,
\newblock \emph{The On-Line Encyclopedia of Integer Sequences},
\newblock \href{https://oeis.org/A174796}{A174796},
Number of admissible sequences of order \(j\); related to \(7x+1\)
problem.

\bibitem{OEIS206788}
OEIS Foundation Inc.,
\newblock \emph{The On-Line Encyclopedia of Integer Sequences},
\newblock \href{https://oeis.org/A206788}{A206788},
Denominators of semiconvergents to \(\log_2 3\).

\bibitem{Wagon85}
S.~Wagon,
\newblock The Collatz problem,
\newblock \emph{Math. Intelligencer} \textbf{7} (1985), 72--76.

\bibitem{WinFerrers}
M.~Winkler,
\newblock Coefficient First-Passage Words for the \(3x+1\) Map:
A Beatty--Ferrers Model and a Power-of-Two Law,
\newblock preprint, 2026.
\newblock \href{https://doi.org/10.13140/RG.2.2.31257.79200}{doi:10.13140/RG.2.2.31257.79200}.

\bibitem{WinBeatty}
M.~Winkler,
\newblock Marked Rotations and Factorization Heights for Dual Beatty
Passage Counts,
\newblock preprint, 2026.
\newblock \href{https://doi.org/10.13140/RG.2.2.22015.57761}{doi:10.13140/RG.2.2.22015.57761}.

\end{thebibliography}

\end{document}
